\subsection{Theory: Big Data at the LHC \quad / \quad \textthai{ทฤษฎี: ข้อมูลขนาดใหญ่ที่ LHC}}
\begin{paracol}{2}
A single collision event at the LHC involves thousands of particles hitting various detector layers. Processing this requires "Fast Inference." Traditionally, we used Cut-based selections (e.g., "keep only if energy > 100 GeV"). However, these simple cuts miss subtle signals like the Higgs Boson or Dark Matter candidates.

Graph Neural Networks (GNNs) are ideal for this task because particle detectors are essentially irregular graphs, not grids. Each sensor hit is a "Node," and the potential path between them is an "Edge."

\switchcolumn

\begin{thai}
เหตุการณ์การชนเพียงครั้งเดียวที่ LHC เกี่ยวข้องกับอนุภาคหลายพันตัวที่พุ่งชนชั้นเครื่องตรวจวัดต่างๆ การประมวลผลสิ่งนี้ต้องการ "การอนุมานที่รวดเร็ว (Fast Inference)" ตามปกติเราใช้การคัดเลือกแบบอิงเกณฑ์ (Cut-based) (เช่น "เก็บไว้เฉพาะเมื่อพลังงาน > 100 GeV") อย่างไรก็ตาม เกณฑ์ง่ายๆ เหล่านี้มักพลาดสัญญาณที่ละเอียดอ่อน เช่น อนุภาคฮิกส์หรือผู้สมัครสสารมืด

Graph Neural Networks (GNNs) เหมาะอย่างยิ่งสำหรับงานนี้ เนื่องจากเครื่องตรวจวัดอนุภาคมีลักษณะเป็นกราฟที่ไม่สม่ำเสมอ ไม่ใช่ตารางพิกเซล เซนเซอร์ที่ถูกชนแต่ละตัวคือ "โหนด (Node)" และเส้นทางที่เป็นไปได้ระหว่างกันคือ "ขอบ (Edge)"
\end{thai}
\end{paracol}
